\chapter{Conclusiones}

Desde el punto de vista técnico, el repositorio documenta un proceso ETL cerrado y reproducible para integrar cuatro fuentes públicas heterogéneas bajo una llave comunal estable. La secuencia de perfilado, homologación, staging, integración, validación y carga final se encuentra respaldada por código fuente, reportes técnicos y artefactos exportados. El resultado es un dataset comunal consistente y una base SQLite consultable, sin necesidad de reinterpretar ni corregir el modelo observado.

Desde el punto de vista analítico, los outputs de Fase 9 muestran que las dimensiones incluidas no se distribuyen homogéneamente entre comunas. Se observan grupos comunales recurrentes en posiciones de rezago relativo, tanto en rankings individuales como en cruces de pobreza con áreas verdes o con IPP por habitante. Estas observaciones deben mantenerse en clave descriptiva, sin extrapolar mecanismos causales.

El valor principal del trabajo no se limita al informe. El CSV final y la base SQLite quedan disponibles como productos reutilizables para consultas, auditoría metodológica, análisis exploratorios posteriores y montaje del informe final en el template externo. En ese sentido, la Fase 10 preparada en este paquete no reemplaza el ETL, sino que lo traduce a materiales documentales trazables y listos para integración editorial.
